\documentclass[aspectratio=169]{beamer}

% Поддержка русского языка и шрифтов
\usepackage[T2A]{fontenc}
\usepackage[utf8]{inputenc}
\usepackage[russian]{babel}
\usepackage{lmodern}
\usepackage{tikz}

% Настройка стильной пастельной палитры
\definecolor{BgPastel}{HTML}{F4F1DE}       % Мягкий кремовый фон (не белый)
\definecolor{TextMain}{HTML}{2B2D42}       % Глубокий тёмно-синий/графитовый для контраста
\definecolor{AccentDrive}{HTML}{E07A5F}    % Драйвовый пастельный терракотовый
\definecolor{AccentTech}{HTML}{3D405B}     % Строгий пастельный индиго
\definecolor{AccentFresh}{HTML}{81B29A}    % Свежий пастельный мятный
\definecolor{AccentLight}{HTML}{F2CC8F}    % Мягкий жёлтый для выделений

% Настройка темы Beamer
\setbeamercolor{background canvas}{bg=BgPastel}
\setbeamercolor{normal text}{fg=TextMain}
\setbeamercolor{frametitle}{fg=AccentDrive}
\setbeamercolor{itemize item}{fg=AccentFresh}
\setbeamercolor{itemize subitem}{fg=AccentDrive}

% Убираем навигационные символы для чистоты дизайна
\setbeamertemplate{navigation symbols}{}

% Настройка заголовков слайдов
\setbeamertemplate{frametitle}{
    \vspace{0.5cm}
    \textbf{\Large\insertframetitle}
}

% Кастомная команда для стильного выделения текста
\newcommand{\hl}[1]{\textcolor{AccentDrive}{\textbf{#1}}}
\newcommand{\tech}[1]{\textcolor{AccentFresh}{\textbf{#1}}}

\begin{document}

% ================= СЛАЙД 1: ТИТУЛЬНЫЙ =================
\begin{frame}
    \begin{tikzpicture}[remember picture,overlay]
        % Декоративные элементы
        \fill[AccentTech] (current page.south west) rectangle ([xshift=4cm,yshift=10cm]current page.south west);
        \fill[AccentDrive] (current page.north east) circle (2.5cm);
        \fill[AccentFresh] (current page.south east) circle (1.5cm);
    \end{tikzpicture}
    
    \vspace{1cm}
    \hspace{3.5cm}
    \begin{minipage}{0.65\textwidth}
        \textcolor{AccentDrive}{\Huge \textbf{IndorTrafficPlan}} \\[0.5em]
        \textcolor{TextMain}{\Large Проектирование организации \\ дорожного движения на максималках} \\[1em]
        \textcolor{AccentTech}{\rule{5cm}{2pt}} \\[1em]
        \textbf{ИндорСофт} | Цифровые решения для дорог
    \end{minipage}
\end{frame}

% ================= СЛАЙД 2: БОЛЬ И РЕШЕНИЕ =================
\begin{frame}{От рутины к чистому инжинирингу}
    \Large
    Создание ПОДД — это сотни знаков, километры разметки и ГОСТы. \\
    Мы сделали так, чтобы \hl{система работала за вас}.
    
    \vspace{0.8cm}
    \normalsize
    \begin{itemize}
        \item \tech{Автоматизация:} Расстановка ТСОДД в пару кликов.
        \item \tech{Два состояния:} Исходное vs Проектное решение — наглядно.
        \item \tech{Умный анализ:} Поиск объектов к переносу и расчет зон видимости.
        \item \tech{Топооснова:} Работа на интернет-картах, DWG и аэрофотосъемке.
    \end{itemize}
\end{frame}

% ================= СЛАЙД 3: ИСКУССТВЕННЫЙ ИНТЕЛЛЕКТ =================
\begin{frame}{ИИ работает, инженер контролирует}
    \begin{columns}
        \begin{column}{0.5\textwidth}
            \vspace{0.5cm}
            \Large \hl{Нейросети в деле}
            \vspace{0.5cm}
            \normalsize
            
            Загружаете видео с видеорегистратора? 
            IndorTrafficPlan \tech{автоматически распознает} стандартные дорожные знаки.
            
            \vspace{0.5cm}
            Меньше ручного ввода — больше времени на принятие проектных решений.
        \end{column}
        \begin{column}{0.5\textwidth}
            \begin{center}
                \begin{tikzpicture}
                    \fill[AccentLight, rounded corners=10pt] (0,0) rectangle (5,4);
                    \node[TextMain, align=center] at (2.5,2) {\textbf{Распознавание} \\ \textbf{по видео} \\ \vspace{0.2cm} \\ \Large };
                \end{tikzpicture}
            \end{center}
        \end{column}
    \end{columns}
\end{frame}

% ================= СЛАЙД 4: ДОКУМЕНТАЦИЯ БЕЗ БОЛИ =================
\begin{frame}{ГОСТы и Приказы — под капотом}
    \Large
    Никакой ручной верстки таблиц. Мастер формирования документации делает всё сам.
    
    \vspace{0.8cm}
    \normalsize
    \begin{itemize}
        \item \hl{Приказ Минтранса РФ № 49:} Все спецификации уже настроены.
        \item \hl{Динамические чертежи:} Меняете проект — чертеж обновляется сам.
        \item \hl{IndorRoadSign:} Встроенный редактор для знаков индивидуального проектирования.
        \item \hl{Экспорт:} Бесшовная передача данных в ГИС IndorRoad и IndorCAD.
    \end{itemize}
\end{frame}

% ================= СЛАЙД 5: О КОМПАНИИ =================
\begin{frame}{ИндорСофт: 20+ лет в дорожной отрасли}
    \begin{columns}
        \begin{column}{0.33\textwidth}
            \begin{center}
                \Huge \hl{4000+} \\
                \normalsize запусков ПО ежедневно
            \end{center}
        \end{column}
        \begin{column}{0.33\textwidth}
            \begin{center}
                \Huge \tech{300к} \\
                \normalsize километров дорог
            \end{center}
        \end{column}
        \begin{column}{0.33\textwidth}
            \begin{center}
                \Huge \hl{100+} \\
                \normalsize крутых специалистов
            \end{center}
        \end{column}
    \end{columns}
    
    \vspace{1cm}
    \begin{center}
        \normalsize Наш софт используют: РОСАВТОДОР, РОСДОРНИИ, АО ВАД и другие лидеры рынка. \\
        \tech{Всё ПО включено в Реестр российского программного обеспечения.}
    \end{center}
\end{frame}

% ================= СЛАЙД 6: CALL TO ACTION =================
\begin{frame}
    \begin{center}
        \Huge \hl{Готовы ускорить проектирование?}
        
        \vspace{1cm}
        \Large \textbf{19 000 р.} \\
        \normalsize Бессрочная лицензия + год обновлений и техподдержки
        
        \vspace{1cm}
        \begin{tikzpicture}
            \fill[AccentFresh, rounded corners=5pt] (0,0) rectangle (6,1.2);
            \node[BgPastel] at (3,0.6) {\textbf{Скачать пробную версию на 30 дней}};
        \end{tikzpicture}
        
        \vspace{0.8cm}
        \normalsize \tech{indorsoft.ru} \quad | \quad 8 800 333-08-05 \quad | \quad sales@indorsoft.ru
    \end{center}
\end{frame}

\end{document}
